\documentclass[11pt,a4paper]{article}
\usepackage[utf8]{inputenc}
\usepackage[spanish]{babel}
\usepackage{amsmath}
\usepackage{amsfonts}
\usepackage{amssymb}
\usepackage{geometry}
\usepackage{enumerate}
\geometry{top=2.5cm, bottom=2.5cm, left=2.5cm, right=2.5cm}

\begin{document}

\begin{center}
    \textbf{\Large SEGUNDO PARCIAL DE FÍSICA BÁSICA} \\[0.2cm]
    \textbf{\large VERSIÓN C} \\[0.4cm]
\end{center}

\noindent \textbf{Estudiante:} \hrulefill \ \textbf{Matrícula:} \underline{\hspace{3cm}} \\[0.2cm]
\textbf{Fecha:} \underline{\hspace{4cm}} \ \textbf{Calificación:} \underline{\hspace{1.5cm}} / 8.0 puntos \\

\hrule
\section*{I. Teoría (Valor: 4.0 puntos - 0.5 puntos cada una)}
\textit{Selecciona la opción correcta en cada una de las siguientes preguntas. No debe haber tachaduras ni enmiendas.}

\begin{enumerate}
    \item \textbf{A diferencia de la masa, el peso de un objeto:}
    \begin{enumerate}[a)]
        \item Es una propiedad intrínseca que no varía con la gravedad local.
        \item Es la fuerza de atracción gravitatoria que ejerce la Tierra sobre él ($w = m \cdot g$).
        \item Se mide en kilogramos en el Sistema Internacional.
        \item Es una cantidad escalar pura sin dirección ni sentido.
    \end{enumerate}

    \item \textbf{En un gráfico de posición en función del tiempo ($x$ vs $t$), la velocidad instantánea en un tiempo determinado equivale a:}
    \begin{enumerate}[a)]
        \item El área total bajo la curva en ese intervalo.
        \item La pendiente de la recta secante que corta la curva.
        \item La pendiente de la recta tangente a la curva en ese punto.
        \item La longitud del vector de aceleración.
    \end{enumerate}

    \item \textbf{En un Movimiento Circular Uniforme (MCU), aunque la rapidez es constante, la velocidad cambia de dirección, lo que implica una aceleración centrípeta dirigida hacia:}
    \begin{enumerate}[a)]
        \item La dirección tangente a la trayectoria circular.
        \item El centro de la trayectoria circular.
        \item El sentido opuesto del movimiento lineal.
        \item El exterior de la circunferencia en dirección radial.
    \end{enumerate}

    \item \textbf{Con respecto a la fuerza de fricción entre dos superficies secas en contacto, se cumple que:}
    \begin{enumerate}[a)]
        \item El coeficiente de fricción cinética ($\mu_c$) es siempre menor que el coeficiente de fricción estática ($\mu_e$).
        \item La fuerza de fricción estática tiene un valor constante e invariable.
        \item La fuerza de fricción es siempre perpendicular a la superficie de contacto.
        \item El coeficiente de fricción depende del área de contacto visible.
    \end{enumerate}

    \item \textbf{La fuerza es una magnitud física de carácter:}
    \begin{enumerate}[a)]
        \item Escalar.
        \item Vectorial.
        \item Adimensional.
        \item Tensorial de segundo orden.
    \end{enumerate}

    \item \textbf{La propiedad física de un cuerpo que mide cuantitativamente su inercia (resistencia al cambio en su estado de movimiento) es la:}
    \begin{enumerate}[a)]
        \item Aceleración.
        \item Fuerza normal.
        \item Masa.
        \item Fricción.
    \end{enumerate}

    \item \textbf{De acuerdo con la Tercera Ley de Newton (Acción-Reacción), si un cuerpo A ejerce una fuerza $\vec{F}_{AB}$ sobre un cuerpo B, entonces:}
    \begin{enumerate}[a)]
        \item El cuerpo B ejerce sobre A una fuerza de igual magnitud, misma dirección y sentido contrario ($\vec{F}_{BA} = -\vec{F}_{AB}$).
        \item El cuerpo B ejerce sobre A una fuerza de mayor magnitud en el mismo sentido.
        \item La fuerza del cuerpo B sobre A es nula porque A inició el contacto.
        \item El cuerpo B ejerce sobre A una fuerza perpendicular a la original.
    \end{enumerate}

    \item \textbf{La Segunda Ley de Newton establece matemáticamente que la aceleración de un cuerpo es:}
    \begin{enumerate}[a)]
        \item Directamente proporcional a su masa e inversamente proporcional a la fuerza neta.
        \item Directamente proporcional a la fuerza neta e inversamente proporcional a su masa.
        \item Independiente de la masa y proporcional al peso.
        \item Proporcional al cuadrado de la fuerza neta aplicada.
    \end{enumerate}

\end{enumerate}

\hrule
\section*{II. Problemas Prácticos (Valor: 4.0 puntos - 1.0 punto cada uno)}
\textit{Resuelve cada uno de los problemas y selecciona la respuesta correcta. Muestra tus cálculos breves.}

\begin{enumerate}
    \item \textbf{Un bloque en una superficie horizontal recibe tres fuerzas paralelas al eje $x$: $F_{1x} = -5.54 \text{ N}$, $F_{2x} = +8.69 \text{ N}$ y $F_{3x} = -3.34 \text{ N}$. Determine el valor de la fuerza neta horizontal:}
    \begin{enumerate}[a)]
        \item $+6.49 \text{ N}$
        \item $-0.19 \text{ N}$
        \item $+0.19 \text{ N}$
        \item $+4.81 \text{ N}$
    \end{enumerate}

    \item \textbf{Un estudiante mide que las aspas de un ventilador realizan exactamente 45 \text{ vueltas} en un intervalo de 2.23 \text{ s}. Calcule la velocidad angular $\omega$ del ventilador en rad/s (utilice $\pi \approx 3.14$):}
    \begin{enumerate}[a)]
        \item $20.2 \text{ rad/s}$
        \item $126.8 \text{ rad/s}$
        \item $106.8 \text{ rad/s}$
        \item $253.6 \text{ rad/s}$
    \end{enumerate}

    \item \textbf{Una cápsula de pasajeros en una rueda de la fortuna de 140 \text{ m} de diámetro gira a rapidez constante y le toma exactamente 35.0 \text{ minutos} (2100 \text{ s}) completar una vuelta entera. Calcule su rapidez lineal:}
    \begin{enumerate}[a)]
        \item $4.000 \text{ m/s}$
        \item $0.104 \text{ m/s}$
        \item $0.209 \text{ m/s}$
        \item $0.418 \text{ m/s}$
    \end{enumerate}

    \item \textbf{Un cuerpo pesa 12.9 \text{ N} en la Tierra ($g = 9.80 \text{ m/s}^2$). ¿Cuánto pesará este mismo cuerpo en la superficie de un asteroide donde la aceleración de la gravedad es de 2.97 \text{ m/s}^2$?}
    \begin{enumerate}[a)]
        \item $4.34 \text{ N}$
        \item $12.9 \text{ N}$
        \item $3.91 \text{ N}$
        \item $8.41 \text{ N}$
    \end{enumerate}

\end{enumerate}

\end{document}
